% (C) 2026 Ralf Stephan, in collaboration with Claude Code.  CC0 1.0.
\chapter{Growth and entropy}\label{ch:entropy}

The number of labelled paths of length \(n\) in a finite carry graph grows either polynomially
or exponentially. Nothing in between occurs, and the rank of Chapter~\ref{ch:reach} decides
which (Theorem~\ref{thm:dichotomy}). Most of the problem families in Chapter~\ref{ch:overview}
ask for the polynomial side, in one of six equivalent forms that this chapter collects in
Theorem~\ref{thm:six-faces}. The exponential side is witnessed by a \emph{horseshoe}: two
closed walks of one length at one vertex (Section~\ref{sec:horseshoe}).

Throughout, \(V\) and \(A\) are finite. Recall \(\mathrm{npaths}(v,n)\), the number of labelled
paths of length \(n\) from \(v\). Two paths differ if they differ in some arc, label or
target, so parallel arcs with different labels count separately.

\section{The dichotomy}

\begin{definition}\label{def:branch}
\(X\) has \emph{polynomial growth}\index{polynomial growth} if \(\mathrm{npaths}(v,n)\le C(n+1)^d\) for some \(C,d\) and
all \(v,n\). A vertex \(v\) \emph{branches}\index{branching} if it has two distinct arcs \(v\to^a t\),
\(v\to^b u\), \((a,t)\neq(b,u)\), whose targets both reach \(v\) again.
\end{definition}

A labelled rank (Definition~\ref{def:rank} with ``successor'' replaced by ``arc'') exists iff no
vertex branches (Theorem~\ref{thm:rank-complete}).

\begin{theorem}[growth dichotomy]\label{thm:dichotomy}\index{growth dichotomy}
\begin{enumerate}[label=(\alph*)]
\item If \(\rho\) is a labelled rank and \(\rho\le R\), then
  \(\mathrm{npaths}(v,n)\le\bigl((|A||V|+1)(n+1)\bigr)^R\).
\item If \(v\) branches, there is \(N>0\) with \(\mathrm{npaths}(v,kN)\ge 2^k\) for all \(k\).
\end{enumerate}
Hence either \(X\) has polynomial growth or some vertex doubles its path count every \(N\)
steps, and not both.
\end{theorem}

\begin{proof}
(a) Fix \(n\le N\) and put \(Y=|A||V|N+1\). We show by induction on \(r\) that a vertex of rank
\(r\) has at most \(Y^r\) paths of length \(n\le N\). From a vertex of rank \(r\), at most one
arc keeps the rank. The other arcs, at most \(|A||V|\) of them, lead to rank \(<r\) and
contribute at most \(Y^{r-1}\) paths each. So along the rank-keeping chain the count grows by
at most \(|A||V|\,Y^{r-1}\) per step, and after \(n\le N\) steps it is at most
\(1+n|A||V|Y^{r-1}\le Y^r\).

(b) Let \(t\leadsto v\) in \(j_1\) steps and \(u\leadsto v\) in \(j_2\) steps. Paths from
\(v\) of length \(m\) inject into paths from \(t\) and from \(u\) of length \(j_1+j_2+1+m\) by
prefixing the return walks. The two first arcs are distinct, so
\(\mathrm{npaths}(v,m+j_1+j_2+2)\ge 2\,\mathrm{npaths}(v,m)\). Take \(N=j_1+j_2+2\).
Exponential beats polynomial, so the two alternatives exclude each other.
\end{proof}

\section{Entropy}

\begin{definition}\label{def:entropy}
The \emph{entropy}\index{entropy} of \(X\) is
\[
  h(X)=\limsup_{n\to\infty}\frac1n\log\sum_{v\in V}\mathrm{npaths}(v,n)\ \in[-\infty,\infty].
\]
\end{definition}

This is the topological entropy of the edge shift of \(X\) \cite[§4.1]{LM95}. It is
\(-\infty\) exactly when the core is empty (Theorem~\ref{thm:core-empty}), and otherwise it is
\(\ge0\). By the Perron--Frobenius theorem it equals \(\log\lambda_A\), where \(\lambda_A\) is the
spectral radius of the adjacency matrix \cite[Thm.~4.3.1, Thm.~4.4.4]{LM95}. This book does not
prove that formula. Nothing in the problem families needs \(\lambda_A\) itself, and the
following brackets give the bounds that the certificates actually check.

\begin{theorem}[eigenvector brackets, {\cite[§5.3]{LM95}}]\label{thm:brackets}\index{eigenvector brackets}
Let \(a,b\in\N\), \(b>0\), and write \((Aw)(v)=\sum_{v\to^c t}w(t)\).
\begin{enumerate}[label=(\alph*)]
\item If \(w\colon V\to\N\) satisfies \(w\ge1\) and \(b\,(Aw)\le a\,w\) pointwise,
  then \(b^n\mathrm{npaths}(v,n)\le a^n w(v)\) and \(h(X)\le\log(a/b)\).
\item If \(a\,w\le b\,(Aw)\) pointwise, \(w\le M\), and \(w(v_0)>0\), then
  \(a^n w(v_0)\le b^n M\,\mathrm{npaths}(v_0,n)\) and \(h(X)\ge\log(a/b)\).
\end{enumerate}
\end{theorem}

\begin{proof}
Both path bounds follow by induction on \(n\) from
\(\mathrm{npaths}(v,n+1)=\sum_{v\to^c t}\mathrm{npaths}(t,n)\). Summing over \(v\) and taking
\(\frac1n\log\), the constants disappear in the limit.
\end{proof}

For example, the golden mean graph (arcs \(0\to^0 0\), \(0\to^1 1\), \(1\to^0 0\)) with
\(w=(8,5)\) satisfies both brackets with \(13/8\) and \(8/5\). So its entropy lies in
\([\log 1.6,\log 1.625]\); in fact it is \(\log\varphi\). The check is a finite computation, and
in Lean it is a single \texttt{decide}. The entropy certificates of the interval problem are
brackets of exactly this kind.

\begin{proposition}\label{prop:entropy-poly}
Polynomial growth implies \(h(X)\le0\). If some vertex satisfies \(\mathrm{npaths}(v,kN)\ge c^k\)
with \(c\ge1\), then \(h(X)\ge(\log c)/N\).
\end{proposition}

\begin{proof}
For every \(\varepsilon>0\), \(C(n+1)^d\le e^{\varepsilon n}\) eventually. The lower bound reads
the \(\limsup\) along the subsequence \(n=kN\).
\end{proof}

\section{Horseshoes}\label{sec:horseshoe}

\begin{definition}\label{def:horseshoe}
A \emph{horseshoe}\index{horseshoe} at \(v\) with petals indexed by a set \(I\) is a family of closed labelled
walks at \(v\), all of one length \(L>0\), that are pairwise distinct as arc sequences.
\end{definition}

For a sequence \(\beta\colon\N\to I\), the \emph{block path} reads petal \(\beta(k)\) in
positions \(kL,\dots,(k+1)L-1\). It is an infinite labelled path from \(v\), because each petal
returns to \(v\), and distinct sequences give distinct paths.

\begin{theorem}\label{thm:horseshoe}
Let \(H\) be a horseshoe at \(v\) with \(|I|\ge1\) petals of length \(L\).
\begin{enumerate}[label=(\alph*)]
\item \(\mathrm{npaths}(v,kL)\ge|I|^k\), so \(h(X)\ge(\log|I|)/L\).
\item If \(|I|\ge2\), some infinite labelled path from \(v\) is not eventually periodic.
\item Conversely, a branching vertex carries a horseshoe with two petals.
\end{enumerate}
\end{theorem}

\begin{proof}
(a) The \(|I|^k\) block paths of \(k\) blocks are distinct paths of length \(kL\). A finite
family of distinct paths of length \(n\) from \(v\) has at most \(\mathrm{npaths}(v,n)\)
members, by induction on \(n\) after sorting the family by first arc. (b) Choose two petals
\(i\neq j\) that first differ at position \(m\), and read petal \(i\) in block \(k\) iff \(k\) is
a power of two. If the path had period \(P\) from \(n_0\) on, then blocks \(k\) and \(k+P\) would
agree at position \(m\) for all \(k\ge n_0\). But for \(k=2^{n_0+P}\), exactly one of \(k\) and
\(k+P\) is a power of two. (c) With \(v\to^a t\leadsto v\) and \(v\to^b u\leadsto v\), the two
loops concatenated in either order have one length and differ in their first arc.
\end{proof}

Part (a) is the horseshoe lower bound of the interval problem, a rose of \(N\) words of length
\(L\). Part (b) is the obstruction of the cell problem: a horseshoe realised by the orbit
dynamics rules out zero entropy for \emph{every} cover of the same system
(Chapter~\ref{ch:covers}).

\section{Zero entropy, six ways}

\begin{theorem}[zero entropy]\label{thm:six-faces}\index{zero entropy}
For a finite carry graph the following are equivalent:
\begin{enumerate}[label=(\arabic*)]
\item a labelled rank exists;
\item no vertex branches;
\item every infinite labelled path is eventually periodic, vertices and labels with one period;
\item \(X\) has polynomial growth;
\item \(h(X)\le0\);
\item every strong component carrying a cycle is a simple cycle: each cyclic vertex has exactly
  one arc whose target returns to it.
\end{enumerate}
\end{theorem}

\begin{proof}
(1)\(\Leftrightarrow\)(2) is Theorem~\ref{thm:rank-complete}. (1)\(\Rightarrow\)(3) is
Theorem~\ref{thm:evper}. (3)\(\Rightarrow\)(2) is Theorem~\ref{thm:horseshoe}(b,c).
(2)\(\Rightarrow\)(4) is Theorem~\ref{thm:dichotomy}(a) with the descendant-count rank.
(4)\(\Rightarrow\)(5) and (5)\(\Rightarrow\)(2) are Proposition~\ref{prop:entropy-poly} with
Theorem~\ref{thm:dichotomy}(b). (2)\(\Leftrightarrow\)(6) unwinds the definitions.
\end{proof}

Applied to the label-forgetting graph \(X^{\circ}\) (one label, an arc \(s\to t\) whenever \(X\)
has one), Theorem~\ref{thm:six-faces} gives the unlabelled form. A rank exists iff no vertex has
two distinct \emph{successors} in its component, iff every infinite walk is eventually periodic,
iff walk counts grow polynomially, iff \(h(X^{\circ})\le0\). For the digit graph of the
coprimality subshift this is the completeness of the digit-problem certificate. A rank
certificate exists iff residue walks grow polynomially, iff every admissible residue walk is
eventually periodic, iff the subshift has zero entropy. The ``only if'' directions were known;
the ``if'' directions are new.

Face (6), with words instead of letters on the arcs, is the phase test of \cite[Thm.~21]{Oda26}.
The word-labelled reading needs subdivision and is postponed to Chapter~\ref{ch:splitting}.

\section*{Notes}
Entropy of edge shifts is \cite[§4.1--4.4]{LM95}. The counting proof of the dichotomy avoids
Perron--Frobenius and is the one the certificates use: the rank is the witness for the
polynomial side and the horseshoe for the exponential side. The irreducible decomposition
\(h(X)=\max_C h(X_C)\) over strong components \cite[Thm.~4.4.4]{LM95} and the period of a
strong component \cite[§10.5]{BJG09} are stated here without Lean proofs. Neither is used later.

\section*{Lean dictionary}
\begin{center}\small
\begin{tabular}{@{}>{\raggedright}p{0.3\textwidth}>{\raggedright\arraybackslash}p{0.64\textwidth}@{}}
\toprule
Text & Lean (\texttt{CarryGraph/Growth.lean}, \texttt{Horseshoe.lean}, \texttt{ZeroEntropy.lean})\\
\midrule
path counts & \lean{arcsFrom}, \lean{npaths\_succ'}, \lean{npaths\_le\_succ},
  \lean{npaths\_mono\_of\_reach}\\
Definition~\ref{def:branch} & \lean{PolyGrowth}, \lean{LNoBranch}\\
Theorem~\ref{thm:dichotomy} & \lean{npaths\_le\_of\_isLRank}, \lean{isLRank\_polyGrowth},
  \lean{pow\_le\_npaths\_of\_branch}, \lean{polyGrowth\_lnoBranch}, \lean{polyGrowth\_or\_exp}\\
Definition~\ref{def:entropy} & \lean{ntot}, \lean{entropy} (Mathlib \lean{expGrowthSup})\\
Theorem~\ref{thm:brackets} & \lean{IsSuperEigen}, \lean{IsSubEigen}, \lean{pow\_mul\_npaths\_le},
  \lean{pow\_mul\_le\_npaths}, \lean{entropy\_le\_of\_superEigen},
  \lean{log\_le\_entropy\_of\_subEigen}\\
Proposition~\ref{prop:entropy-poly} & \lean{expGrowthSup\_poly\_le},
  \lean{entropy\_le\_zero\_of\_polyGrowth}, \lean{log\_div\_le\_entropy},
  \lean{entropy\_le\_zero\_iff}, \lean{entropy\_pos\_iff}\\
Definition~\ref{def:horseshoe} & \lean{Horseshoe}, \lean{Horseshoe.blockS}, \lean{Horseshoe.blockL},
  \lean{Horseshoe.isIPath}, \lean{Horseshoe.block\_injective}\\
Theorem~\ref{thm:horseshoe} & \lean{card\_le\_npaths}, \lean{Horseshoe.pow\_le\_npaths},
  \lean{Horseshoe.log\_card\_div\_le\_entropy}, \lean{Horseshoe.exists\_not\_evPer},
  \lean{exists\_horseshoe\_of\_branch}\\
Theorem~\ref{thm:six-faces} & \lean{zeroEntropy\_tfae}, \lean{AllEvPer}, \lean{SimpleCycles},
  \lean{lnoBranch\_iff\_simpleCycles}\\
unlabelled form & \lean{noBranch\_tfae}, \lean{lnoBranch\_forget\_iff}, \lean{isLRank\_forget\_iff}\\
digit problem & \lean{CarryBridge.certOK\_iff\_polyGrowth}, \lean{CarryBridge.certOK\_iff\_allEvPer},
  \lean{CarryBridge.certOK\_iff\_entropy}, \lean{CarryBridge.certOK\_iff\_entropy\_labelled}\\
witness checking & \lean{decIsRank}, \lean{decIsLRank}, \lean{decLNoBranch},
  \lean{decIsSuperEigen}, \lean{decIsSubEigen}\\
\bottomrule
\end{tabular}
\end{center}
